\section{Las tres medias de frontera}\label{app:boundary}
En \(\Gamma\) se ha aislado cada cuadrado, pero sus dos pares diagonales pueden estar disponibles
a la vez. Para guardar el signo y contar una sola vez su interacción, empezamos por la
corrección de un par y después reunimos los dos pares en el cuadrado. Aquí \(p_e\) designa
la ocupación estacionaria de \emph{torres} en la máscara binaria.
Para un par diagonal disponible \(P=\{e,f\}\), sea \(\pi^{\rm ex}_P\) su probabilidad conjunta en
una colocación uniforme de torres. Sea \(\pi^T_P=\exp[L_T(M_P)-L_T(M)]\), donde \(M_P\) elimina sus
cuatro extremos de fila y columna. Con \(\mathsf G(e,f)\) la respuesta cruzada del punto estacionario
de torres y \(\chi_e\) su corrección relativa de una casilla, definimos
\[
 \pi_{1,P}=p_ep_f(1-\mathsf G(e,f)+\chi_e+\chi_f),\qquad
 X_P=p_ep_f(p_e+p_f)\mathsf G(e,f).
\]
Para un par solo la corrección del cuadrado es
\(\log(1-\pi^{\rm ex}_P)-\log(1-\pi_{1,P})-X_P+O(j^{-5})\).
Las probabilidades son \(O(j^{-2})\) y su diferencia es \(O(j^{-4})\); al integrar la derivada
logarítmica el error de linealización es \(O(j^{-6})\). El cuadrado completo da las correcciones
de sus dos pares y un término adicional \(-3\prod_{s\in\square}p_s+O(j^{-5})\).
Hay a lo sumo \(j^2(j-1)\) pares, pues cada casilla tiene a lo sumo \(2(j-1)\) vecinas y cada
par tiene dos extremos. Así, la siguiente expresión \(B_n(b)\) difiere de \(n\mathbb E_\mu[\Gamma;F_b]\)
en \(O(n/j^2)=o_b(1)\):
\begin{align*}
 B_n(b)={}&-n\mathbb E_\mu[\textstyle\sum_P(\pi^{\rm ex}_P-\pi^T_P);F_b]\\
           &-n\mathbb E_\mu[\textstyle\sum_P(\pi^T_P-\pi_{1,P}+X_P);F_b]
            -3n\mathbb E_\mu[\textstyle\sum_\square\prod_{s\in\square}p_s;F_b].
\end{align*}
La inserción de \(\pi^T_P\) en esta descomposición tiene dos funciones. La primera media mide
el error del predictor de torres en el par. La segunda conserva el error de inserción
para un par marcado. Para promediarlo hay que saber cuáles de los otros dos vértices siguen
disponibles. Sean \(u_f\) las vacancias finales de las dos familias diagonales y
\(U_f(b)=u_f+b(1-u_f)\). Sean \(v_b=U_D(b)U_A(b)\), \(p_b\) su escalado
bistocástico~\cite{sk} y
\(W_b^0=p_b/v_b\). Un cuadrado de coordenadas mínimas \((r,c)\) y lado \(h\) se integra sobre
\(\Omega=\{0<h<1,\ 0<r,c<1-h\}\).
Con \(d_s=U_D(b,r-c+sh)\), \(a_s=U_A(b,r+c+sh)\), la disponibilidad de su par orientado
tiene factor \(B_o\). Los otros dos vértices tienen indicadores \(I_l I_u,I_l I_v\).
Sus cuatro probabilidades conjuntas y el factor del par marcado son los siguientes.
\[
\begin{array}{c|cccc}
o&B_o&l&u&v\\\hline
 D&d_0a_0a_2&a_1&d_{-1}&d_{+1}\\
 A&a_1d_{-1}d_{+1}&d_0&a_0&a_2
\end{array}
\qquad
\begin{array}{c|c}
s&q_{o,s}\\\hline
00&1-l+l(1-u)(1-v)\\
10&lu(1-v)\\
01&l(1-u)v\\
11&luv
\end{array}
\]
La submedida de cada estado es \(b^4 B_o q_{o,s}\,dr\,dc\,dh\). Las cuatro etiquetas vivas son
dos filas y dos columnas. Una vacancia diagonal repetida conserva su mismo indicador.
Por orientación hay \(\sum_{h=1}^{n-1}(n-h)^2\) posiciones.
La normalización de Riemann aporta \(n^3\), con cuatro factores \(j^{-1}\) y un \(n\) exterior.
Por tanto,
\(n^4 b^4/j^4\to1\). Al insertar los dos pesos marcados resulta
\(w_{o,b}=B_oW_b^0(e)W_b^0(f)=p_b(e)p_b(f)/U_o(b)\).
Se conservan incluso los estados de masa nula sin dividir por \(q_{o,s}\).

Para la primera media, al integrar \(p_b\) sobre cada línea obtenemos su masa \(m_f\). Sean
\[
 W_b=\sum_{f=D,A}\int\frac{m_f^2}{2U_f},\qquad
 \gamma_b=\frac{m_D}{U_D}+\frac{m_A}{U_A},\qquad
 \theta_b(r)=\int p_b(r,c)\gamma_b(r,c)\,dc .
\]
Eliminar los cuatro extremos en la expansión de torres da el término de tamaño \(W_b/6\), el
de normalización \(-2F(b)W_b\) y la variación de las medidas de fila y columna \(-dF_2(b)\).
Aquí \(dF_2(b)\) es la derivada del funcional de torres a núcleo fijo al multiplicar ambas
medidas uniformes por \(1+\varepsilon(2W_b-\theta_b)\), con la simetría de columna.
La dirección tiene masa cero por \(\int\theta_b=2W_b\). Resulta
\(P_T=\lim_{b\downarrow0}[(1/6-2F(b))W_b-dF_2(b)]\).
Para la tercera media, el cuadrado usa seis etiquetas diagonales distintas. Dos aparecen dos
veces entre los ocho factores de vacancia de \(\prod_s W_b^0(s)\). Al multiplicar por las seis
vacancias distintas y tomar \(p_0=\nu_1\) queda exactamente
\[
 -S=-3\int_\Omega
       \frac{\prod_{s\in\square}\nu_1(s)}
            {u_D(r-c)\,u_A(r+c+h)}\,dr\,dc\,dh .
\]
Así quedan identificadas la primera y la tercera medias. La segunda conserva dos datos de la
máscara en el cuadrado marcado, los indicadores de los cruces. El apéndice~\ref{app:weights}
calcula el error del par con esos dos indicadores fijos.
